%%%%%%%%%%%%%%%%%%%%%%%%%%%%%%%%%%%%%%%%%%%%%%%%%%%%%%%%%%%%%%%%%%%%%%%%%%%%%%%%%%%%%%%%%%%%%%%%%%%%%%%%%%%%%%%%%%%%%%%%%%%%%%%%%%%%%%%%%%%%%%%%%%%%%%%%%%%
% This is just an example/guide for you to refer to when producing your supplementary material for your Frontiers article.                                 %
%%%%%%%%%%%%%%%%%%%%%%%%%%%%%%%%%%%%%%%%%%%%%%%%%%%%%%%%%%%%%%%%%%%%%%%%%%%%%%%%%%%%%%%%%%%%%%%%%%%%%%%%%%%%%%%%%%%%%%%%%%%%%%%%%%%%%%%%%%%%%%%%%%%%%%%%%%%

%%% Version 2.5 Generated 2022/06/14 %%%
%%% You will need to have the following packages installed: datetime, fmtcount, etoolbox, fcprefix, which are normally included in WinEdt. %%%
%%% In http://www.ctan.org/ you can find the packages and how to install them, if necessary. %%%
%%%  NB logo1.jpg is required in the path in order to correctly compile front page header %%%

\documentclass[utf8]{frontiers_suppmat} % for all articles
% T1 font encoding -- without it, \textendash/\textemdash (used by the
% literatecard style's literate= mapping) fall back to OT1-encoded glyphs
% whose PDF ToUnicode mapping is wrong, corrupting en/em dash in extracted
% PDF text (confirmed via a real pdflatex+pdftotext round trip,
% WI-CHARTER-SUPPLEMENT-RENDERER-DASH-FIX).
\usepackage[T1]{fontenc}
% Outline (scalable) fonts, not the class's default bitmap Computer Modern --
% pdftotext silently drops dashes and quotes extracted from bitmap-font text
% (confirmed in the 2026-10-26 charter supplement pipeline audit), which
% would make any PDF-text regression check unreliable.
\usepackage{lmodern}
\usepackage{url,hyperref,lineno,microtype}
\usepackage[onehalfspacing]{setspace}
% graphicx (already loaded by frontiers_suppmat.cls for the front-page logo,
% loaded again explicitly here so this file doesn't depend on that class
% internal) -- provides \scalebox, used below to fix the en/em dash glyphs
% WI-CHARTER-SUPPLEMENT-LISTING-DASH-GLYPH-FIX found still visually wrong.
\usepackage{graphicx}


\usepackage{framed}
\usepackage{xcolor}
\newenvironment{Shaded}{\begin{shaded}}{\end{shaded}}
\definecolor{shadecolor}{rgb}{0.95, 0.95, 0.95}
\providecommand{\tightlist}{%
  \setlength{\itemsep}{0pt}\setlength{\parskip}{0pt}}

\usepackage{listings}

\lstset{
  basicstyle=\ttfamily\small,
  breaklines=true,
  frame=single,
  columns=fullflexible
}

\lstdefinelanguage{YAML}{
  keywords={true,false,null},
  sensitive=false,
  comment=[l]{\#},
  morestring=[b]",
  morestring=[b]'
}

\lstset{
  language=YAML,
  basicstyle=\ttfamily\small,
  breaklines=true,
  frame=single,
  columns=fullflexible,
  showstringspaces=false,
  keywordstyle=\color{blue},
  commentstyle=\color{darkgray},
  stringstyle=\color{teal}
}

\newcommand{\yamlkey}[1]{\texttt{\color{blue}#1}}

% WI-CHARTER-SUPPLEMENT-LISTING-DASH-GLYPH-FIX: \textendash/\textemdash
% inside this monospace (\ttfamily) context are the FONT's own correctly
% T1-encoded en/em dash glyphs -- pdftotext already extracts them as the
% right Unicode codepoint (U+2013/U+2014) -- but Latin Modern Mono draws
% those glyphs visually as short as a hyphen (confirmed directly: an
% isolated pdflatex+pdftoppm comparison against the same characters in
% body prose, where they render correctly long). This is a genuine
% monospace-font design limitation, not a mapping bug -- the same short
% rendering was independently confirmed in TeX Gyre Cursor, a different
% monospace face with its own distinct T1 endash/emdash glyphs. DejaVu
% Sans Mono renders a visually longer mark, but its TeX Live T1 encoding
% vector substitutes /hyphen for /endash at that slot (confirmed via the
% font's own .enc file), so it extracts as a plain hyphen -- not usable.
% \dashstretch keeps the real \textendash/\textemdash glyph (so the
% correct ToUnicode mapping is untouched -- \scalebox only scales how an
% already-selected glyph is drawn, not which glyph or character it is)
% and stretches it horizontally so it visually reads as a true dash
% instead of a hyphen. Verified via a real pdflatex+pdftoppm round trip,
% inside an actual framed, line-wrapped lstlisting, with no overlap or
% alignment artifacts.
\newcommand{\dashstretch}[1]{\scalebox{1.8}[1]{#1}}

\lstdefinestyle{literatecard}{
  basicstyle=\ttfamily\small,
  breaklines=true,
  frame=single,
  columns=fullflexible,
  showstringspaces=false,
  keepspaces=true,
  literate=
    {–}{{\dashstretch{\textendash}}}1
    {—}{{\dashstretch{\textemdash}}}1
    {‑}{{-}}1
    {“}{{\textquotedblleft}}1
    {”}{{\textquotedblright}}1
    {’}{{\textquoteright}}1
}

% Leave a blank line between paragraphs instead of using \\

\begin{document}
\onecolumn
\firstpage{1}

\title[Supplementary Material]{{\helveticaitalic{Supplementary Material}}}

\maketitle

% \section{Supplementary Data}

% Supplementary Material should be uploaded separately on submission. Please include any supplementary data, figures and/or tables. All supplementary files are deposited to FigShare for permanent storage and receive a DOI.

%Supplementary material is not typeset so please ensure that all information is clearly presented, the appropriate caption is included in the file and not in the manuscript, and that the style conforms to the rest of the article. To avoid discrepancies between the published article and the supplementary material, please do not add the title, author list, affiliations or correspondence in the supplementary files.

% \section{Supplementary Tables and Figures}

% For more information on Supplementary Material and for details on the different file types accepted, please see \href{http://home.frontiersin.org/about/author-guidelines#SupplementaryMaterial}{the Supplementary Material section} of the Author Guidelines.

% Figures, tables, and images will be published under a Creative Commons CC-BY licence and permission must be obtained for use of copyrighted material from other sources (including re-published/adapted/modified/partial figures and images from the internet). It is the responsibility of the authors to acquire the licenses, to follow any citation instructions requested by third-party rights holders, and cover any supplementary charges.

%% Figures, tables, and images will be published under a Creative Commons CC-BY licence and permission must be obtained for use of copyrighted material from other sources (including re-published/adapted/modified/partial figures and images from the internet). It is the responsibility of the authors to acquire the licenses, to follow any citation instructions requested by third-party rights holders, and cover any supplementary charges.
This section contains normative cards from the Prosocial Robot Navigation Charter available at
\url{https://github.com/xenotaur/prosoc/}, formatted for display in LaTeX by the Pandoc package available at \url{https://pandoc.org/} and lightly edited for display. These cards represent a snapshot of the system in its current state of development, and are presented to show the expressiveness and scope of the system, not to present fully-audited normative card content suitable for deployment.

@@CHARTER@@

\section{Social Navigation Scenario Cards}
@@FRONTAL_APPROACH@@

@@SINGLE_FILE_HALLWAY@@

@@MOVABLE_OBSTRUCTION@@

\section{Social Navigation Task Cards}

@@NAVIGATE_POINT_TO_POINT@@

@@DELIVER_OBJECT@@

\section{Social Navigation Context Cards}
@@ROUTINE_DELIVERY@@

@@GUIDANCE_DOCENT@@

@@HIGH_URGENCY@@

%\subsection{Figures}

%%% There is no need for adding the file termination, as long as you indicate where the file is saved. In the examples below the files (logo1.eps and logos.eps) are in the Frontiers LaTeX folder
%%% If using *.tif files convert them to .jpg or .png
%%%  NB logo1.eps is required in the path in order to correctly compile front page header %%%

%\begin{figure}[htbp]
%\begin{center}
%\includegraphics[width=9cm]{logo1}% This is a *.eps file
%\end{center}
%\caption{ Enter the caption for your figure here.  Repeat as  necessary for each of your figures}\label{fig:1}
%\end{figure}

%\setcounter{figure}{2}
%\setcounter{subfigure}{0}
%\begin{subfigure}
%\setcounter{figure}{2}
%\setcounter{subfigure}{0}
%    \centering
%    \begin{minipage}[b]{0.5\textwidth}
%        \includegraphics[width=\linewidth]{logo1.eps}
%        \caption{This is Subfigure 1.}
%        \label{fig:Subfigure 1}
%    \end{minipage}
%
%\setcounter{figure}{2}
%\setcounter{subfigure}{1}
%    \begin{minipage}[b]{0.5\textwidth}
%        \includegraphics[width=\linewidth]{logo2.eps}
%        \caption{This is Subfigure 2.}
%        \label{fig:Subfigure 2}
%    \end{minipage}%\

%\setcounter{figure}{2}
%\setcounter{subfigure}{-1}
%    \caption{Enter the caption for your subfigure here. %\textbf{(A)} This is the caption for Subfigure 1. %\textbf{(B)} This is the caption for Subfigure 2.}
%    \label{fig: subfigures}
%\end{subfigure}

%%% If you don't add the figures in the LaTeX files, please upload them when submitting the article.
%%% Frontiers will add the figures at the end of the provisional pdf automatically
%%% The use of LaTeX coding to draw Diagrams/Figures/Structures should be avoided. They should be external callouts including graphics.


%\bibliographystyle{Frontiers-Harvard} %  Many Frontiers journals use the Harvard referencing system (Author-date), to find the style and resources for the journal you are submitting to see \href{https://zendesk.frontiersin.org/hc/en-us/articles/360017860337-Frontiers-Reference-Styles-by-Journal}{Frontiers Reference Styles by Journal}. For Humanities and Social Sciences articles please include page numbers in the in-text citations
%\bibliographystyle{Frontiers-Vancouver} % Many Frontiers journals use the numbered referencing system, to find the style and resources for the journal you are submitting to see \href{https://zendesk.frontiersin.org/hc/en-us/articles/360017860337-Frontiers-Reference-Styles-by-Journal}{Frontiers Reference Styles by Journal}.

%\bibliography{test}

\end{document}
